\section{相关理论与技术基础}

\subsection{无人机集群仿真技术概述}

无人机集群系统拥有一些核心特征，这些特征共同形成了它与传统单机系统不同的核心竞争力：

其一，分布性方面，集群内多架无人机在物理空间上是分散部署的状态，没有中心化控制节点，每架无人机都搭载着独立的感知单元，像视觉、雷达、GPS这类，还有计算单元。借助swarm\_test.py脚本能够模拟多机分布式感知数据的传输与处理，以此验证分布式架构的稳定性。

其二，协同性方面，集群依靠多机间的信息共享，比如位置、姿态、任务进度、环境障碍等信息，来达成任务协作，能够依据全局任务目标动态分配子任务，举例来说，部分无人机负责区域侦察，部分负责精准作业，部分负责通信中继，通过协同调度实现"1+1>2"这样的整体效能。

其三，实时性。无人机集群在指令下发、信息交互、状态反馈这些环节，要满足毫秒级的严格时延约束。特别是在高速移动、动态的环境中，要是时延过高，就会产生轨迹偏差、避障失效等问题。为了保障实时性，需要对集群通信链路、计算单元加以优化。

其四，鲁棒性。集群具备容错和自修复能力，要是部分单机出现传感器故障、通信中断或者动力失效等状况，系统能够依靠协同算法重新分配任务、规划轨迹，从而保证整体任务不会中断。

无人机集群的核心技术包括集群控制、路径规划、协同感知、避障决策、通信组网、任务分配、资源调度等方面。这些技术在实际应用之前要经过大量仿真测试，以确保安全性、可靠性、有效性。

\subsection{无人机集群仿真平台分类与对比}

无人机集群仿真平台按技术架构可分为两大类：基于游戏引擎的高保真仿真平台和基于动力学模型的轻量化仿真平台。

\begin{figure}[H]
    \centering
    \includegraphics[width=0.9\textwidth]{figures/image3.png}
    \caption{无人机集群仿真平台分类框架图}
    \label{fig:classification}
\end{figure}

这两类平台对比如表\ref{tab:platform}所示：

\begin{table}[H]
    \centering
    \caption{无人机集群仿真主流平台对比表}
    \label{tab:platform}
    \begin{tabular}{|p{2.5cm}|p{2.5cm}|p{1cm}|p{1cm}|p{1.5cm}|p{1cm}|p{2.5cm}|}
        \hline
        平台类型 & 代表平台 & 物理精度 & 视觉效果 & 多机支持 & 开源性 & 适用场景 \\
        \hline
        高保真视觉 & AirSim & 高 & 极高 & 差（未优化） & 开源 & 视觉、控制、规划 \\
        \hline
        高保真视觉 & Unity/UE自定义 & 中高 & 极高 & 自定义 & 商业 & 高端演示、虚拟仿真 \\
        \hline
        轻量化 & ROS/Gazebo & 中 & 低 & 中 & 开源 & 算法快速验证 \\
        \hline
        轻量化 & PX4 SITL & 中高 & 低 & 中 & 开源 & 飞控算法验证 \\
        \hline
        轻量化 & MATLAB & 高 & 低 & 中 & 商业 & 理论仿真、分析 \\
        \hline
    \end{tabular}
\end{table}

注：高保真视觉平台侧重渲染与物理真实性，轻量化平台侧重算法快速验证；多机支持能力为大规模集群仿真的关键指标。

AirSim凭借高物理保真、高视觉保真、ROS兼容、多载体支持、开源免费等优势\cite{fang2025,li2020,xiao2020,yang2024}，成为本研究的最优选择。

大规模无人机集群仿真的性能优劣，直接决定算法验证有效性、控制决策可靠性以及场景还原真实性，需从实时性、资源效率、时序准确性、运行鲁棒性等多个维度综合评判。本研究结合无人机集群实时仿真需求，构建包含实时性指标、资源占用指标、时序精度指标、系统鲁棒性指标在内的完整性能评价体系，具体核心指标包括：帧率（FPS）、仿真时延、CPU占用率、线程上下文切换频率、定时精度、稳定性、可扩展性。

\subsection{AirSim仿真平台原理与架构}

AirSim（Aerial Informatics and Robotics Simulation）是微软研究院2017年发布的开源仿真平台，最初面向自动驾驶汽车，后续扩展支持多旋翼无人机、固定翼无人机、无人车、机器人等多种载体。

AirSim基于物理引擎与游戏引擎构建，提供高精度动力学模型、传感器仿真（摄像头、激光雷达、IMU、GPS、磁力计）、天气系统（光照、雨、雪、雾）、多机协同接口、Python/C++ API、ROS兼容，可用于感知、控制、规划、SLAM、强化学习、集群协同等算法验证。

\begin{figure}[H]
    \centering
    \includegraphics[width=0.7\textwidth]{figures/image4.png}
    \caption{无人机渲染加载模块验证情况图}
    \label{fig:render}
\end{figure}

\begin{figure}[H]
    \centering
    \includegraphics[width=0.7\textwidth]{figures/image5.png}
    \caption{无人机动态模块加载模块验证情况图}
    \label{fig:load}
\end{figure}

AirSim采用分层架构设计，分为引擎层、核心仿真层、接口层、应用层。

\begin{figure}[H]
    \centering
    \includegraphics[width=0.85\textwidth]{figures/image6.png}
    \caption{AirSim平台整体架构图}
    \label{fig:architecture}
\end{figure}

AirSim的主循环按照固定频率运行，承担着无人机物理更新、状态更新、指令处理、数据反馈等工作。每一台无人机都对应着一个独立的代理对象，这个代理对象包括状态、控制、传感器、通信等多个模块。AirSim针对多无人机集群仿真的原生默认实现，运用了轻量化架构设计，具体的实现方式如下：

其一，采用单机单线程模型，为每一台无人机实例创建独立线程，以此达成多机并行仿真。

其二，线程内部借助无限循环结合sleep(0)来完成时序同步与空闲等待。

其三，由统一的主循环驱动全局物理更新、环境交互、状态刷新。

其四，通过标准API接口来完成多机指令分发、状态采集、数据交互。

该实现在小规模集群场景中有着实现简洁、部署便捷等优点，然而在面向大规模集群仿真时存在性能缺陷：线程数量会随着无人机规模线性增长，容易超出CPU物理核心的承载能力，sleep(0)高频触发线程调度，使得用户态与内核态切换开销上升，在超线程CPU环境里，自旋等待机制失效，容易引发调度风暴，多线程资源竞争加剧，上下文切换频繁，系统有效吞吐量大幅降低。上述这些问题构成了大规模集群仿真的关键性能瓶颈，是本研究重点要解决的核心问题。

\subsection{高性能并发与线程调度理论}

线程是进程内的独立执行单元\cite{chen2023}，是操作系统调度的最小单位。同一进程内的线程共享进程资源，切换开销远小于进程。

线程包含五种基本状态：新建、就绪、运行、阻塞、终止。

\begin{figure}[H]
    \centering
    \includegraphics[width=0.5\textwidth]{figures/image7.png}
    \caption{线程状态转换示意图}
    \label{fig:thread}
\end{figure}

线程调度指的是操作系统挑选就绪线程来运行的过程，这里面涵盖了抢占式调度、时间片轮转、优先级调度等策略。线程调度会引发上下文切换开销，像寄存器保存、页表更新、缓存刷新等都在其中，频繁进行切换会严重降低系统性能。

各类机制性能对比如表\ref{tab:waiting}所示：

\begin{table}[H]
    \centering
    \caption{线程等待机制性能特征对比表}
    \label{tab:waiting}
    \begin{tabular}{|p{2.5cm}|p{2cm}|p{1.5cm}|p{2cm}|p{3cm}|}
        \hline
        同步机制 & 上下文切换开销 & CPU资源占用 & 定时/响应精度 & 典型适用场景 \\
        \hline
        sleep\_for & 中等 & 低 & 中等 & 低实时性需求、长时间等待 \\
        \hline
        yield & 较高 & 中等 & 中低 & 短时间等待、高并发调度 \\
        \hline
        sleep(0) & 极高 & 高 & 低 & 仅简易同步，高并发不推荐 \\
        \hline
        纯自旋 & 无 & 极高 & 高 & 极短等待、强实时场景 \\
        \hline
        条件变量 & 中等 & 低 & 中等 & 线程间协同、事件触发 \\
        \hline
    \end{tabular}
\end{table}

操作系统分为用户态与内核态。应用程序执行系统调用（如sleep、yield、文件操作、网络操作）时，会触发用户态→内核态→用户态切换。每次切换需要保存上下文、切换页表、刷新缓存，产生固定开销。

在AirSim的默认实现里，有大量线程频繁去执行sleep(0)、yield这类系统调用，使得内核态占比超过了60\%，CPU资源大量被浪费在切换上面，这便是性能下降的关键原因。

自旋等待是一种忙等式同步机制，线程在等待的时候会持续执行循环检测，不会主动释放CPU控制权。阻塞睡眠是一种被动式等待机制，线程借助系统调用进入阻塞状态，主动释放CPU资源，由操作系统内核进行调度唤醒。

AirSim原生多机仿真默认采用sleep(0)混合等待策略，在工程实现上存在明显不合理性：该方式既无法达到自旋等待的高精度与低时延特性，又会频繁触发系统调用与线程调度，带来极高的上下文切换开销与内核态负载，直接导致大规模集群仿真性能急剧下降。

\subsection{性能监测与分析工具}

Windows Performance Analyzer（WPA）是Windows操作系统自带的系统级高性能可视化分析工具，依托Windows Event Trace（ETW）内核事件追踪机制，可对系统运行时行为进行高精度、低侵入式的数据采集与可视化分析，是定位多线程调度、CPU资源占用、系统开销等性能瓶颈的关键工具。

WPA可捕获并解析的核心指标包括：全局与进程级CPU占用率、线程状态转换、上下文切换频率、进程与线程生命周期、函数调用栈深度与热点函数分布。

perf是Linux内核原生集成的开源性能事件监测与分析工具，基于内核性能计数器与跟踪子系统实现，具备低侵入、高精度、覆盖维度广等特点，广泛用于多线程调度、CPU开销、内存访问、缓存命中率等关键性能问题定位。

本研究借助Python/C++混合编程达成了一套具有低侵入性、实时性、可扩展性的自定义性能采集与统计模块。此模块能够在仿真运行进程中实时采集关键指标，为优化效果比对、瓶颈定位、结论验证提供直接的数据支持。

\subsection{本章小结}

本章系统阐述无人机集群仿真、AirSim平台架构、线程调度、同步等待、内核态开销、性能监测工具等相关理论，为后续性能瓶颈定位、优化方案设计、实验测试评估提供坚实理论支撑与技术基础。
